% =====================================================================
%  Chapter 6 — Principle: Resilience (Fail Gently, Never Silently)
%  Map: research/chapter-map.md (Chapter 6)
%  Shape: idea → origin → cost → what it opens up. Ideas, not how-to.
%  Rule: the author is never quoted in body text; short quotes only in
%  footnotes. Full statements live in interviews/.
% =====================================================================
\chaptersubtitle{Fail Gently, Never Silently}
\chapter{Resilience}
\label{ch:failsoft}

Silica's motto has two halves. The first, secure by default at compile time,
belongs to Chapters~\ref{ch:danger} and~\ref{ch:refuses}. The second is about
what happens when something goes wrong anyway. The program should fail
gently---one small part, not the whole system---and it should never fail
silently.\footnote{The motto, from \src{silica/README.md}: ``Secure by default
at compile time --- fail soft, never fail silent.''}

\section{Let it fall, help it up}

The idea comes from Erlang, and it is the most counter-intuitive in this book.
Rather than surrounding every operation with defensive code, let the part that
fails simply stop, and have a supervisor start it again from a known good
state. Reliability comes from containment and recovery, not from trying to
foresee every mistake.

Silica's beginner book explains it with an elephant mother. She does not stop
her calf from ever stumbling. When it falls, she gets it back on its feet,
fresh.\footnote{\src{silica/docs/learn-programming.md}, lines 27--29 and 1339.}
Silica takes the idea out of Erlang's virtual machine and builds it into
programs that run directly on the hardware, and it insists that every program
arrange for failures to be noticed before anything else begins. Barney had
learned the hard way, building SwErl, that supervision is nearly impossible
without control over how a failing actor reports its failure
(Chapter~\ref{ch:actors}).

% TODO-INTERVIEW: Where did the elephant-mother image come from?

\section{Guard only what can be ruined}

Letting things fail does not mean never checking. Silica's guidance names just
two reasons to guard against an error: when the failure would corrupt data, and
when it would happen so often that recovering every time would slow the system
down. Everywhere else, the advice is to write no check at all and let the
supervisor do its work. When something does go wrong---a request for an item
past the end of a list, for instance---the faulty operation is never carried
out, so nothing is damaged, and only the part that asked for it stops.

% TODO-VERIFY: The two-criteria wording is agent-recorded in .claude-memory;
%   the beginner book has a published version (learn-programming.md ~1331).

\section{Never silent}

The second half of the motto is held just as firmly. A failure is always
reported, and there is no way to switch the report off. Silica regards a
failure that leaves no trace as a fault in the language's own runtime. Failures
a program expects are handled as ordinary results; failures it does not expect
stop the part concerned and are announced.

\section{Whoever sets a limit owns it}

A related idea runs through Silica's treatment of resources. The language
itself imposes few arbitrary ceilings: an actor's memory grows as needed, and
queues of waiting messages are not capped. The reasoning is that limits belong
to whoever chose them. A program that sets its own sizes, or starts a thousand
risky workers, knows its numbers, and handling the consequences is that
program's job, not a hidden rule's.

\section{When the ground itself is broken}

Letting a part fail assumes the ground beneath it is still sound, and foreign
code written in C is the hardest test of that assumption. Silica meets it with
the same idea it uses everywhere else. Foreign code runs only inside workers set
aside for it, and when that code crashes, the runtime catches the fault, ends
that one worker, throws away the scratch memory it shared with the C code, and
hands the failure to the worker's supervisor, which restarts it like any other
actor. The rest of the application carries on.\footnote{The design names the
dangerous actor as ``the desired failure unit'':
\src{silica/compiler/design_documents/dangerous_ffi_security_model.md}, lines
217 and 243--244.}

Silica is candid about where that promise ends. C can damage memory without
crashing at all, or damage the runtime itself, and no supervisor can repair
what cannot be detected. Only when a fault strikes outside those guarded
workers, or the runtime can no longer trust its own state, does the whole
program stop---and even then it says why. For libraries that cannot be trusted
at all, the design points to running them in a separate process altogether.

% TODO-INTERVIEW: How did you arrive at the guarded dangerous actor as the
%   unit of failure for foreign code? What did it take to make that work?

\section{A motto that changed}

When the project went public in April 2026, its motto said in effect that bad
security practice should always be a compiler error. Within days it became the
two-part motto above, and it later gained the words at compile time. The change
gave the running program, not only the compiler, a part in keeping its
promises.

% TODO-INTERVIEW: What changed in your thinking between the two mottos?

\section{What it costs}

Letting things fail asks programmers to unlearn habits of defensive coding,
and asks designers to think about supervision from the first day. Refusing to
fail silently means facing every failure, even those it might be more
comfortable not to know about.

\section{What it opens up}

\begin{itemize}
  \item \textbf{Systems that keep running.} Devices in the field, services
        that cannot go down, and machines no one can reach to restart can heal
        themselves part by part.
  \item \textbf{Simpler code.} Without layers of defensive checks, the code
        that remains says what the program is for.
  \item \textbf{Failures that teach.} Because every failure is reported, a
        system's problems become visible and fixable instead of rumoured.
  \item \textbf{Reliability on small machines.} Recovery by supervision
        works on a microcontroller as well as on a server.
  \item \textbf{Old libraries without old fragility.} A crash in a trusted
        C library costs one worker and a restart, not the whole
        application.
\end{itemize}

% TODO-INTERVIEW: Where would you most like to see this kind of resilience
%   that it doesn't exist today?
